% ════════════════════════════════════════════════════════════════
% CAPÍTULO 4 -- INTEGRAIS DUPLOS E MODELAÇÃO COMPUTACIONAL
% Ficheiro para \include{} no template principal.
%
% NOTA PARA O TEMPLATE PRINCIPAL: Este capítulo utiliza código
% MATLAB/Octave. Adicionar ao preâmbulo do ficheiro principal:
%
%   \usepackage{listings}
%   \lstset{
%       language=Matlab,
%       basicstyle=\ttfamily\footnotesize,
%       keywordstyle=\color{blue}\bfseries,
%       commentstyle=\color{teal},
%       stringstyle=\color{red},
%       numbers=left,
%       numberstyle=\tiny\color{gray},
%       stepnumber=1,
%       breaklines=true,
%       frame=single,
%       backgroundcolor=\color{gray!5},
%       captionpos=b,
%       showstringspaces=false
%   }
% ════════════════════════════════════════════════════════════════

\chapter{Integrais Duplos e Modelação Computacional}

% ────────────────────────────────────────────────────────
\section{Fundamentos Teóricos}

% ────────────────────────────────────────────────────────
\subsection{O Conceito de Integral Duplo}

Enquanto um integral simples calcula a área debaixo de uma curva num gráfico 2D, um integral duplo calcula o volume abaixo de uma superfície num espaço 3D. A notação padrão é dada por:

\begin{defbox}[title={Integral Duplo}]
\[
    V = \iint_R f(x,y) \,\dd A
\]
onde $R$ representa a região planar base (o domínio) e $f(x,y)$ a função que dita a altura da superfície em cada ponto $(x,y)$.
\end{defbox}

% ────────────────────────────────────────────────────────
\subsection{Definição Formal e Somas de Riemann}

A construção rigorosa de um integral duplo baseia-se na expansão das Somas de Riemann. Seja $f(x,y)$ uma função definida numa região planar fechada e limitada $R$. Se dividirmos $R$ numa malha de pequenos retângulos de área $\Delta A$, e escolhermos um ponto de amostragem $(x_i^*, y_j^*)$ dentro de cada sub-retângulo, o volume aproximado do sólido abaixo da superfície é a soma dos volumes de vários paralelepípedos:

\begin{defbox}[title={Soma de Riemann Dupla}]
\[
    V \approx \sum_{i=1}^{n} \sum_{j=1}^{m} f(x_i^*, y_j^*) \Delta A
\]
O integral duplo é definido como o limite desta soma quando a área dos sub-retângulos tende para zero (ou seja, quando a malha se torna infinitamente fina):
\[
    \iint_R f(x,y) \,\dd A = \lim_{\Delta A \to 0} \sum_{i=1}^{n} \sum_{j=1}^{m} f(x_i^*, y_j^*) \Delta A
\]
\end{defbox}

\begin{notebox}
Computacionalmente, os métodos de integração numérica (como a regra dos trapézios bidimensional usada nos projetos) são aproximações algorítmicas diretas desta mesma Soma de Riemann.
\end{notebox}

% ────────────────────────────────────────────────────────
\subsection{Teorema de Fubini, Integração Iterada e Tipos de Regiões}

O cálculo direto pelo limite da soma de Riemann é impraticável. O Teorema de Fubini é a ponte que permite calcular integrais duplos transformando-os em dois integrais simples sucessivos. Este teorema resolve o problema afirmando que, se a função $f(x,y)$ for contínua na região $R$, o integral duplo pode ser avaliado através de integrais iterados unidimensionais. A complexidade do cálculo depende da geometria da região $R$:

\begin{thmbox}[title={Teorema de Fubini e Integração Iterada}]
\begin{itemize}
    \item \textbf{Regiões do Tipo I (Simples na vertical):} A região é limitada à esquerda e à direita por constantes $x=a$ e $x=b$, e em cima e em baixo por funções contínuas $y=g_1(x)$ e $y=g_2(x)$. Integra-se primeiro em ordem a $y$:
    \[
        \iint_R f(x,y) \,\dd A = \int_{a}^{b} \left[ \int_{g_1(x)}^{g_2(x)} f(x,y) \,\dd y \right] \dd x
    \]

    \item \textbf{Regiões do Tipo II (Simples na horizontal):} A região é limitada em cima e em baixo por constantes $y=c$ e $y=d$, e lateralmente por funções $x=h_1(y)$ e $x=h_2(y)$. Integra-se primeiro em ordem a $x$:
    \[
        \iint_R f(x,y) \,\dd A = \int_{c}^{d} \left[ \int_{h_1(y)}^{h_2(y)} f(x,y) \,\dd x \right] \dd y
    \]
\end{itemize}
\end{thmbox}

\begin{tipbox}
Dependendo da geometria da região $R$, a capacidade de alterar a ordem de integração (de Tipo I --- $\dd x\,\dd y$ para Tipo II --- $\dd y\,\dd x$) é uma competência analítica vital, pois muitas vezes um integral iterado é impossível de resolver analiticamente numa ordem, mas trivial na outra, podendo este processo simplificar drasticamente a resolução do problema.
\end{tipbox}

% ────────────────────────────────────────────────────────
\subsection{Transformação para Coordenadas Polares}

Quando a região de integração possui simetria circular ou a função integranda contém termos como $x^2 + y^2$, o uso de coordenadas cartesianas $(x,y)$ gera fronteiras de integração complexas com raízes quadradas. Nestes casos, aplica-se uma transformação geométrica para o sistema de coordenadas polares $(r, \theta)$, onde:
\[
    x = r \cos\theta \quad \text{e} \quad y = r \sin\theta
\]

\begin{thmbox}[title={Mudança para Coordenadas Polares e Jacobiano}]
Ao deformarmos o espaço bidimensional para um novo sistema de coordenadas, a área dos elementos infinitesimais deforma-se. O fator de escala dessa deformação é dado pelo valor absoluto do determinante Jacobiano da transformação. Para coordenadas polares, o Jacobiano $J$ calcula-se como:
\[
    J = \begin{vmatrix}
    \dfrac{\partial x}{\partial r} & \dfrac{\partial x}{\partial \theta} \\[6pt]
    \dfrac{\partial y}{\partial r} & \dfrac{\partial y}{\partial \theta}
    \end{vmatrix}
    = \begin{vmatrix}
    \cos\theta & -r\sin\theta \\
    \sin\theta & r\cos\theta
    \end{vmatrix}
    = r(\cos^2\theta + \sin^2\theta) = r
\]
Portanto, o elemento de área diferencial muda de $\dd x\,\dd y$ para $r \,\dd r \,\dd\theta$. O teorema geral da mudança de variáveis estabelece que:
\[
    \iint_R f(x,y) \,\dd x\,\dd y = \iint_{S} f(r\cos\theta, r\sin\theta) \, r \,\dd r\,\dd\theta
\]
onde $S$ é a representação da região $R$ no novo plano $r\theta$.

Em suma:
\[
    \dd A = \dd x \, \dd y = r \, \dd r \, \dd\theta
\]
\end{thmbox}

% ────────────────────────────────────────────────────────
\section{Modelação Computacional e Projetos Interdisciplinares}

A verdadeira avaliação das competências ocorre quando os conceitos analíticos são traduzidos em algoritmos. Computadores resolvem integrais através de métodos numéricos (somatórios em malhas discretas). Abaixo detalham-se os três projetos de aplicação prática.

% ────────────────────────────────────────────────────────
\subsection{Projeto 1: Simulação de Difusão de Calor}

A energia térmica numa superfície planar acumula-se e difunde-se ao longo do tempo. Para calcular o calor total acumulado, integra-se a função de temperatura sobre a área 2D.

\begin{lstlisting}[caption={Simulação de Difusão de Calor em MATLAB/Octave}]
% 1. Configuracao do Dominio (Regiao Planar R)
nx = 50; ny = 50; dx = 0.1; dy = 0.1;
T = zeros(nx, ny);
T(20:30, 20:30) = 100; % Fonte de calor inicial

% 2. Simulacao da Difusao (Diferencas Finitas)
alpha = 0.2;
for k = 1:50
    T_nova = T;
    for i = 2:nx-1
        for j = 2:ny-1
            T_nova(i,j) = T(i,j) + alpha * (T(i+1,j) + T(i-1,j) + T(i,j+1) + T(i,j-1) - 4*T(i,j));
        end
    end
    T = T_nova;
end

% 3. Calculo do Integral Duplo Numerico (Energia Total)
calor_total = sum(sum(T)) * (dx * dy);

% 4. Visualizacao 3D da Superficie
surf(T); colormap('hot');
title(sprintf('Energia Total: %.2f', calor_total));
\end{lstlisting}

% ────────────────────────────────────────────────────────
\subsection{Projeto 2: Análise do Movimento Solar e Vetores 3D}

Para estimar a radiação total capturada por um painel, modela-se o painel como uma região $R$ no espaço e calcula-se o ângulo de incidência do vetor solar. A energia total diária exige integração dupla sobre a área do painel iterada ao longo do tempo.

\begin{lstlisting}[caption={Análise de Incidência Solar em MATLAB/Octave}]
x = 0:0.1:2; y = 0:0.1:1.5; % Dimensoes do painel
[X, Y] = meshgrid(x, y);
normal_painel = [0, 0, 1];
energia_total = 0;

for t = 0:1:12
    theta = (pi/12) * t;
    vetor_sol = [cos(theta), 0, sin(theta)];

    % Intensidade via produto interno (Vetores 3D)
    intensidade = max(0, dot(normal_painel, vetor_sol));

    % Funcao densidade de energia f(x,y) (ex: lado sombreado)
    Sombra = 1 - 0.1 * X;
    Energia_XY = intensidade * Sombra;

    % Integral Duplo iterado (Metodo dos Trapezios)
    energia_instante = trapz(y, trapz(x, Energia_XY, 2));
    energia_total = energia_total + energia_instante;
end
\end{lstlisting}

% ────────────────────────────────────────────────────────
\subsection{Projeto 3: Redes Neurais e Descida de Gradiente}

Ao treinar uma rede neural, otimiza-se uma função de custo $L(w_1, w_2)$ que representa uma superfície no espaço dos parâmetros. A compreensão das derivadas parciais e da topologia de funções de várias variáveis é crucial para o algoritmo de \textit{Gradient Descent}.

\begin{lstlisting}[caption={Otimização em Superfície de Erro (Redes Neurais)}]
% 1. Espaco de Parametros (Pesos w1, w2)
[w1, w2] = meshgrid(-5:0.2:5, -5:0.2:5);
Custo = w1.^2 + 2 * w2.^2; % Funcao de Custo (Superficie 3D)

% 2. Otimizacao via Descida de Gradiente
peso = [4, 4]; % Ponto inicial
taxa = 0.1;
hist_pesos = zeros(21, 2); hist_pesos(1,:) = peso;

for i = 1:20
    grad = [2*peso(1), 4*peso(2)]; % Derivadas parciais
    peso = peso - taxa * grad;
    hist_pesos(i+1,:) = peso;
end

% 3. Curvas de Nivel no Plano 2D
contour(w1, w2, Custo, 20); hold on;
plot(hist_pesos(:,1), hist_pesos(:,2), 'r-o', 'LineWidth', 2);
title('Descida de Gradiente na Superficie de Custo');
\end{lstlisting}